
\section{Pubblicazione di documentazione (sul repo di documentazione e nel sito web)}

Riferirsi alla procedura alla sezione \ref{mainsec:contrib} per quanto riguarda l'utilizzo di git.

\newpage

\subsection{Creazione del documento}
Per pubblicare qualsiasi documento relativo al team git HAPPENS, per prima cosa è necessario trovare (o creare)
la cartella di destinazione nel repo \textit{documents} (e.g. \textit{verbali/esterni/yyyy-mm-dd} o \textit{verbali/interni/yyyy-mm-dd}).


\paragraph{Intestazione e template}
Creare nella cartella un file \textit{main.tex} e inserire la direttiva \texttt{\textbackslash documentclass\{gh\}}.
Questa direttiva importerà in automatico l'intestazione e il template di git HAPPENS.

Combinare in modo opportuno questi comandi per riempire i campi comuni del file (omettere i comandi non necessari):

\begin{center}
    \renewcommand{\arraystretch}{1.3}
    \begin{tabularx}{\textwidth}{@{} l X @{}}
        \toprule
        \textbf{Comando}                                & \textbf{Descrizione}                                           \\
        \midrule
        \texttt{\textbackslash tipodoc\{\dots\}}        & Tipo o titolo del documento, riportato anche nel piè di pagina \\
        \texttt{\textbackslash datadoc\{\dots\}}        & Data del documento                                             \\
        \texttt{\textbackslash modalita\{\dots\}}       & Modalità di svolgimento (es.\ Online - Google Meet)            \\
        \texttt{\textbackslash orario\{\dots\}}         & Orario di inizio e fine                                        \\
        \texttt{\textbackslash redattori\{\dots\}}      & Redattori del documento                                        \\
        \texttt{\textbackslash verificatori\{\dots\}}   & Verificatori del documento                                     \\
        \texttt{\textbackslash amministratore\{\dots\}} & Amministratore del gruppo                                      \\
        \texttt{\textbackslash responsabile\{\dots\}}   & Responsabile, riportato anche nel piè della prima pagina       \\
        \texttt{\textbackslash destinatari\{\dots\}}    & Destinatari del documento                                      \\
        \bottomrule
    \end{tabularx}
\end{center}

Per inserire più valori in ogni campo, utilizzare il separatore "\textbackslash \textbackslash" (e.g. "F. Lissandrin \textbackslash \textbackslash{ } G. Gambarro").

\paragraph{Corpo del documento}
Inserire le direttive di corpo:

\begin{verbatim}
    \begin{document}

    \makeheader


    \end{document}
\end{verbatim}

Per ogni capitolo del file, creare un file in \textit{capitoli/xx\_nome\_capitolo.tex} (dove \textit{xx} è il numero progressivo di capitolo).
Richiamare ogni capitolo inserendo l'istruzione \texttt{ \textbackslash input\{capitoli/xx\_nome\_capitolo.tex\} } dopo l'istruzione \textit{makeheader}.

\subsubsection{Altre convenzioni file}
\begin{itemize}
    \item Se il documento deve contenere immagini, porle in una cartella \textbf{"assets"} allo stesso livello del file \textit{main.tex}
    \item (TBA)
\end{itemize}

\subsection{Pubblicazione del documento}
\begin{enumerate}
    \item (Opzionale) Creare un file \textit{metadata.json} (vedi \ref{docpub:metadata_json})
    \item Facendo riferimento alla procedura \ref{mainsec:contrib}, inviare il documento per la revisione.
\end{enumerate}
Un workflow di GitHub Actions si occuperà di controllare la build, e al merge della \textit{PR} un altro workflow effettuerà la pubblicazione sul sito web.


\subsubsection{File metadata.json}\label{docpub:metadata_json}
Il file \textit{metadata.json} definisce alcuni campi per migliorare la visualizzazione del file nel sito web.

Un file metadata è composto come segue:
\begin{verbatim}
    {
        "title": "Titolo documento (e.g. 'Verbale esterno del 20 ottobre 2026')",
        "description": "Descrizione documento (e.g. 'Primo incontro con Sanmarco Informatica: ...')
    }
\end{verbatim}

L'assenza di un file metadata non causa errori ma peggiora la visualizzazione del file nel sito (titolo dedotto dal nome file e descr. assente).

\subsection{Strumenti di supporto}\label{docpub_tech}

\begin{enumerate}
    \item Estensione \href{https://marketplace.visualstudio.com/items?itemName=James-Yu.latex-workshop}{Latex Workshop} per VS Code: consente di avviare una preview live dei documenti su cui si lavora. Prevede dei requisiti ulteriori da installare, riferirsi alla pagina dell'estensione.
\end{enumerate}